% GENERATED FILE. Edit canonical lessons, then run npm run generate:books.
% canonical-source-hash: fnv1a64:78e63cdc1ded6c18
% canonical-lessons: GU-C278-javu-joie, GU-C278-etle, GU-C278-tethi-etle, GU-R278-advice-and-result

\chapter{Should, must, and so}
\label{ch:gu-should-must-and-so}
\hlchaptermodality{278}

\begin{chapteropening}
\textbf{By the end of this chapter:} \emph{I can say what someone should do and has to do (tamāre javũ joīe, māre javũ paḍśe), and join a reason to its result with eṭle or tethī.}

Complete GU-R278-advice-and-result: give advice, say what has to be done, and say what follows.
\end{chapteropening}

\section[\texorpdfstring{\gu{તમારે} \gu{જવું} \gu{જોઈએ} (tamāre javũ joīe)}{તમારે જવું જોઈએ (tamāre javũ joīe)}]{\gu{તમારે} \gu{જવું} \gu{જોઈએ} (tamāre javũ joīe) — should, and have to}

\label{lesson:GU-C278-javu-joie}



\begin{quote}
Say \textquotedblleft{}because of the rain\textquotedblright{}. (\emph{varsādne līdhe}.) Then \textquotedblleft{}Arun drank tea\textquotedblright{}. (\emph{Aruṇe cā pīdhī}.) Keep that \textbf{-\gu{એ}}.
\end{quote}

\begin{grammarlens}[title={Grammar Lens: the one who should takes -\gu{એ}}]
\begin{quote}
\textbf{\gu{તમારે} \gu{દાક્તર} \gu{પાસે} \gu{જવું} \gu{જોઈએ}.}
\end{quote}

\begin{quote}
\emph{tamāre dāktar pāse javũ joīe.} — \textquotedblleft{}You should go to the doctor.\textquotedblright{}
\end{quote}

The verb's name in \textbf{-\gu{વું}}, then \textbf{\gu{જોઈએ}} (\emph{joīe}), which stays as it is. The one who should takes an \textbf{-\gu{એ}} form:

\noindent
\begin{tabularx}{\linewidth}{@{}>{\raggedright\arraybackslash}X>{\raggedright\arraybackslash}X@{}}
\toprule
\textbf{who} & \textbf{should go} \\
\midrule
I & \textbf{\gu{મારે} \gu{જવું} \gu{જોઈએ}} \\
you, to a friend & \textbf{\gu{તારે} \gu{જવું} \gu{જોઈએ}} \\
he, she & \textbf{\gu{તેણે} \gu{જવું} \gu{જોઈએ}} \\
we & \textbf{\gu{અમારે} \gu{જવું} \gu{જોઈએ}} \\
you, with respect & \textbf{\gu{તમારે} \gu{જવું} \gu{જોઈએ}} \\
Arun, Mira & \textbf{\gu{અરુણે}}, \textbf{\gu{મીરાએ}} \\
\bottomrule
\end{tabularx}

With an object, the \textbf{-\gu{વું}} follows the thing, as a past does: \textbf{\gu{તમારે} \gu{દવા} \gu{લેવી} \gu{જોઈએ}} (\emph{tamāre davā levī joīe}), feminine for \textbf{\gu{દવા}}; \textbf{\gu{મારે} \gu{પત્ર} \gu{લખવો} \gu{જોઈએ}} (\emph{māre patra lakhvo joīe}), masculine for \textbf{\gu{પત્ર}}.
\end{grammarlens}

\subsection*{You'll want to know: have to}
\begin{quote}
\textbf{\gu{મારે} \gu{કાલે} \gu{વહેલા} \gu{જવું} \gu{પડશે}.}
\end{quote}

\begin{quote}
\emph{māre kāle vahelā javũ paḍśe.} — \textquotedblleft{}I'll have to go early tomorrow.\textquotedblright{}
\end{quote}

Put \textbf{\gu{પડવું}} in place of \textbf{\gu{જોઈએ}} and \textquotedblleft{}should\textquotedblright{} becomes \textquotedblleft{}have to\textquotedblright{}: \textbf{\gu{પડે} \gu{છે}}, as a rule; \textbf{\gu{પડશે}}, later; \textbf{\gu{પડ્યું}}, it was so. \textbf{\gu{મારે} \gu{રોજ} \gu{બસમાં} \gu{જવું} \gu{પડે} \gu{છે}} (\emph{māre roj basmā̃ javũ paḍe che}); \textbf{\gu{ગઈકાલે} \gu{મારે} \gu{ચાલીને} \gu{જવું} \gu{પડ્યું}} (\emph{gaīkāle māre cālīne javũ paḍyũ}), \textquotedblleft{}yesterday I had to walk\textquotedblright{}.

\subsection*{Your turn}
Say these aloud:

\begin{itemize}
\raggedright
  \item \textquotedblleft{}He should rest\textquotedblright{} — \emph{teṇe ārām karvo joīe}
  \item \textquotedblleft{}Yesterday I had to walk\textquotedblright{} — \emph{gaīkāle māre cālīne javũ paḍyũ}
\end{itemize}

\begin{tcolorbox}[breakable,colback=teal!4,colframe=teal!35!black,title={Before you move on}]
\textquotedblleft{}You should go\textquotedblright{}? (\textbf{\gu{તમારે} \gu{જવું} \gu{જોઈએ}}.) \textquotedblleft{}I have to go every day\textquotedblright{}? (\textbf{\gu{મારે} \gu{રોજ} \gu{જવું} \gu{પડે} \gu{છે}}.) Why \textbf{\gu{લેવી}} in \textbf{\gu{તમારે} \gu{દવા} \gu{લેવી} \gu{જોઈએ}}? (It follows \textbf{\gu{દવા}}.)
\end{tcolorbox}
\section[\texorpdfstring{\gu{એટલે} (eṭle)}{એટલે (eṭle)}]{\gu{એટલે} (eṭle) — so; that is why}

\label{lesson:GU-C278-etle}



\begin{quote}
Before the new one: say \textquotedblleft{}it was raining\textquotedblright{}. (\emph{varsād paḍto hato}.) Then \textquotedblleft{}that much\textquotedblright{}. (\textbf{\gu{એટલું}}, \emph{eṭlũ}.)
\end{quote}

\subsection*{You'll want to know: \gu{એટલે}}
\textbf{\gu{એટલે}} — \emph{eṭle} — \textquotedblleft{}so; that is why\textquotedblright{}. The reason comes first, then \textbf{\gu{એટલે}}, then what followed from it.

\begin{quote}
\textbf{\gu{વરસાદ} \gu{પડતો} \gu{હતો}, \gu{એટલે} \gu{અમે} \gu{ઘરે} \gu{રહ્યા}.}
\end{quote}

\begin{quote}
\emph{varsād paḍto hato, eṭle ame ghare rahyā.} — \textquotedblleft{}It was raining, so we stayed at home.\textquotedblright{}
\end{quote}

\begin{quote}
\textbf{\gu{હું} \gu{થાકી} \gu{ગઈ} \gu{હતી}, \gu{એટલે} \gu{હું} \gu{ઘરે} \gu{રહી}.}
\end{quote}

\begin{quote}
\emph{hũ thākī gaī hatī, eṭle hũ ghare rahī.} — \textquotedblleft{}I was tired, so I stayed at home.\textquotedblright{}
\end{quote}

\begin{cousinweb}
\textbf{\gu{એટલે}} (\emph{eṭle}) is probably built from \textbf{\gu{એટલું}}, \textquotedblleft{}that much\textquotedblright{}, with the \textbf{-\gu{ે}} of \textbf{\gu{સવારે}}: \textquotedblleft{}to that point\textquotedblright{}, and so \textquotedblleft{}for that reason\textquotedblright{}. CDIAL lists Gujarati \emph{eṭlũ} under Sanskrit \emph{iyattaka}, \textquotedblleft{}so little\textquotedblright{} (CDIAL 1589); it has no row for \textbf{\gu{એટલે}} itself.
\end{cousinweb}

\subsection*{Your turn}
Say these aloud:

\begin{itemize}
\raggedright
  \item \emph{eṭle}
  \item \textquotedblleft{}It was raining, so we stayed at home\textquotedblright{} — \emph{varsād paḍto hato, eṭle ame ghare rahyā}
  \item \textquotedblleft{}The train was late, so I came late\textquotedblright{} — \emph{ṭren moḍī hatī, eṭle hũ moḍo āvyo}
\end{itemize}

\begin{tcolorbox}[breakable,colback=teal!4,colframe=teal!35!black,title={Before you move on}]
What does \textbf{\gu{એટલે}} mean? (\textquotedblleft{}So; that is why\textquotedblright{}.) Which comes first, the reason or the result? (The reason.)
\end{tcolorbox}
\section[\texorpdfstring{\gu{તેથી} · \gu{એટલે} (tethī · eṭle)}{તેથી · એટલે (tethī · eṭle)}]{\gu{તેથી}, \gu{એટલે} (tethī, eṭle) — so, therefore}

\label{lesson:GU-C278-tethi-etle}



\begin{quote}
Say \textquotedblleft{}so\textquotedblright{} in the everyday way. (\textbf{\gu{એટલે}}.) Then \textquotedblleft{}therefore\textquotedblright{}, from A2. (\textbf{\gu{તેથી}}.)
\end{quote}

\begin{grammarlens}[title={Grammar Lens: reason, then result}]
\begin{quote}
\textbf{\gu{તમને} \gu{તાવ} \gu{છે}, \gu{એટલે} \gu{તમારે} \gu{આરામ} \gu{કરવો} \gu{જોઈએ}.}
\end{quote}

\begin{quote}
\emph{tamne tāv che, eṭle tamāre ārām karvo joīe.} — \textquotedblleft{}You have a fever, so you should rest.\textquotedblright{}
\end{quote}

\textbf{\gu{એટલે}} is the everyday \textquotedblleft{}so\textquotedblright{}. \textbf{\gu{તેથી}} (\emph{tethī}), \textquotedblleft{}therefore\textquotedblright{}, says the same and suits a letter or a notice: \textbf{\gu{ટ્રેન} \gu{મોડી} \gu{હતી}, \gu{તેથી} \gu{અમે} \gu{મોડા} \gu{પહોંચ્યા}} (\emph{ṭren moḍī hatī, tethī ame moḍā pahõcyā}).
\end{grammarlens}

\subsection*{You'll want to know: turning a reason round}
\textbf{\gu{કારણ} \gu{કે}} puts the reason after; \textbf{\gu{એટલે}} puts it before. The same two facts, either way:

\noindent
\begin{tabularx}{\linewidth}{@{}>{\raggedright\arraybackslash}X>{\raggedright\arraybackslash}X@{}}
\toprule
\textbf{result first} & \textbf{reason first} \\
\midrule
\textbf{\gu{હું} \gu{મોડો} \gu{આવ્યો}, \gu{કારણ} \gu{કે} \gu{ટ્રેન} \gu{મોડી} \gu{હતી}.} & \textbf{\gu{ટ્રેન} \gu{મોડી} \gu{હતી}, \gu{એટલે} \gu{હું} \gu{મોડો} \gu{આવ્યો}.} \\
\textbf{\gu{અમે} \gu{ઘરે} \gu{રહ્યા}, \gu{કારણ} \gu{કે} \gu{વરસાદ} \gu{પડતો} \gu{હતો}.} & \textbf{\gu{વરસાદ} \gu{પડતો} \gu{હતો}, \gu{એટલે} \gu{અમે} \gu{ઘરે} \gu{રહ્યા}.} \\
\bottomrule
\end{tabularx}

\subsection*{Your turn}
Say these aloud:

\begin{itemize}
\raggedright
  \item \textquotedblleft{}You have a fever, so you should rest\textquotedblright{} — \emph{tamne tāv che, eṭle tamāre ārām karvo joīe}
  \item \textquotedblleft{}The bus was late, so I had to walk\textquotedblright{} — \emph{bas moḍī hatī, eṭle māre cālīne javũ paḍyũ}
\end{itemize}

\begin{tcolorbox}[breakable,colback=teal!4,colframe=teal!35!black,title={Before you move on}]
\textquotedblleft{}So\textquotedblright{}, in everyday speech? (\textbf{\gu{એટલે}}.) In a letter? (\textbf{\gu{તેથી}}.) Turn \textbf{\gu{અમે} \gu{ઘરે} \gu{રહ્યા}, \gu{કારણ} \gu{કે} \gu{વરસાદ} \gu{પડતો} \gu{હતો}} round. (\textbf{\gu{વરસાદ} \gu{પડતો} \gu{હતો}, \gu{એટલે} \gu{અમે} \gu{ઘરે} \gu{રહ્યા}}.)
\end{tcolorbox}
\section[(review)]{Review — should, must, and so}

\label{lesson:GU-R278-advice-and-result}



\begin{quote}
Say \textquotedblleft{}I think it will rain\textquotedblright{}. (\emph{mane lāge che ke varsād paḍśe}.) Then \textquotedblleft{}because of the rain\textquotedblright{}. (\emph{varsādne līdhe}.)
\end{quote}

\subsection*{Your turn}
Say these aloud:

\begin{itemize}
\raggedright
  \item You should go to the doctor. You should take the medicine.
  \item He should rest. I should write a letter.
  \item I have to go by bus every day. Tomorrow I'll have to go early.
  \item What do you think? I think you have a fever, so you should rest.
  \item The train was late, so I came late. Now turn it round, with kāraṇ ke.
\end{itemize}

\begin{tcolorbox}[breakable,colback=teal!4,colframe=teal!35!black,title={Before you move on}]
\textquotedblleft{}I should go\textquotedblright{}? (\textbf{\gu{મારે} \gu{જવું} \gu{જોઈએ}}.) \textquotedblleft{}I had to go\textquotedblright{}? (\textbf{\gu{મારે} \gu{જવું} \gu{પડ્યું}}.) Turn \textbf{\gu{હું} \gu{મોડો} \gu{આવ્યો}, \gu{કારણ} \gu{કે} \gu{ટ્રેન} \gu{મોડી} \gu{હતી}} round with \textbf{\gu{એટલે}}. (\textbf{\gu{ટ્રેન} \gu{મોડી} \gu{હતી}, \gu{એટલે} \gu{હું} \gu{મોડો} \gu{આવ્યો}}.)
\end{tcolorbox}
